\documentclass{article}
\usepackage[UTF8]{ctex}
\usepackage{amsmath,amssymb,amsfonts}
\usepackage{enumitem}
\usepackage{geometry}
\geometry{a4paper, margin=1in}
% === 错题本排版优化（不参与题目标识） ===
\setlength{\parindent}{0pt}
\setlength{\parskip}{0.35em plus 0.10em minus 0.08em}
\setlength{\emergencystretch}{3em}
\setlist{itemsep=0.25em, topsep=0.40em, parsep=0pt}
\allowdisplaybreaks[2]
% === 错题本排版优化结束 ===
\begin{document}
\section*{浙江大学2025年高等代数真题}
\section*{95．浙江大学2025年高等代数试卷}
\begin{itemize}

\subsection{题目1}
  \item[1．] 设实系数首一多项式 \(f(x)\) 在实数上的值处处恒正．证明：存在实系数多项式 \(g(x)\) 和 \(h(x)\) ，使得
\[
f(x)=g^{2}(x)+h^{2}(x) .
\]

\subsection{题目2}
  \item[2．] 记 \(a_{i i}=i^{3}+3 i^{2}+2 i, a_{i j}=\left\{\begin{array}{l}3(i-j)+1, i<j ; \\ 3(i+j)+2, i>j\end{array}\right.\) 证明： 2025 阶行列式 \(\operatorname{det}\left(a_{i j}\right)\) 为 3 的整数倍。
\end{itemize}

\subsection{题目3}
3．设 \(A=\left(a_{i j}\right)_{(n-1) \times n}, b_{i}=(-1)^{i} c_{i}, i=1,2, \cdots, n\) ，其中 \(c_{i}\) 为 \(A\) 去掉第 \(i\) 列所得矩阵的行列式。证明\(\eta=\left(\begin{array}{c}b_{1} \\ \vdots \\ b_{n}\end{array}\right)\) 为 \(A\left(\begin{array}{c}x_{1} \\ \vdots \\ x_{n}\end{array}\right)=0\) 的基础解系的充要条件是 \(\operatorname{rank}(A)=n-1\) 。

\subsection{题目4}
4．设 \(A\) 为 \(n\) 阶矩阵，证明：存在对称矩阵 \(S\) 和可逆矩阵 \(R\) ，满足 \(A=S R\) 。

\subsection{题目5}
5．设 \(A\) 为 \(n\) 阶实矩阵，若 \(A+A^{\mathrm{T}}\) 正定，证明：\(A\) 可逆．

\subsection{题目6}
6．设 4 维复空间 \(V\) 上的可逆线性变换 \(\varphi\) 满足 \(\operatorname{tr}\left(\varphi^{i}\right)=0, i=1,2,3\) 。证明：\(\varphi\) 可对角化。

\begin{itemize}

\subsection{题目7}
\end{itemize}

\subsection{题目7}
原题边界“7．”：求 \(P(\lambda), Q(\lambda)\) 满足 \(P(\lambda) A Q(\lambda)=\operatorname{diag}\left\{1, \lambda-1,(\lambda-1)^{2}\right\}\) ，其中
\[
A=\left(\begin{array}{ccc}
\lambda+1 & 2 & -6 \\
1 & \lambda & -3 \\
1 & 1 & \lambda-4
\end{array}\right)
\]

\subsection{题目8}

\subsection{题目8}
原题边界“8．”：证明：欧氏空间长度相同的向量，存在正交变换，将一个向量映射成另一个向量．

\subsection{题目9}
9．设 \(\mathcal{T}: \mathbb{R}^{4} \rightarrow \mathbb{R}^{3}\) 满足

\[
\mathcal{P}\left(\begin{array}{l}
x_{1} \\
x_{2} \\
x_{3} \\
x_{4}
\end{array}\right)=\left(\begin{array}{cccc}
1 & -2 & 3 & -1 \\
3 & 1 & 1 & -3 \\
1 & 5 & -5 & -1
\end{array}\right)\left(\begin{array}{l}
x_{1} \\
x_{2} \\
x_{3} \\
x_{4}
\end{array}\right)
\]

记 \(V, W\) 分别为 \(\mathbb{R}^{4}, \mathbb{R}^{3}\) 的对偶空间。定义 \(\mathcal{T}^{*}: W \rightarrow V\) 满足 \(\mathcal{T}^{*}(f)=f \circ \mathcal{T}\) 。证明： \(\mathcal{T}^{*}\) 为线性映射，并求 \(\operatorname{Ker} \mathcal{T}^{*}\) 与 \(\operatorname{Im} \mathcal{T}^{*}\) 。

\subsection{题目10}
10．\(A, B\) 是 \(n\) 阶复矩阵，\(A-E\) 是幂零矩阵，若存在正整数 \(k\) 使得 \(A^{k} B=B A^{k}\) ，证明：\(A B=B A\) 。
\end{document}
